\paragraph{Theorem31.} There exists an efficient randomized algorithm such that for every$k$ and any$\varepsilon>0$, outputs with probability$1-2^{-\Omega(k)}$ an$\varepsilon$-balanced linear code$C:\mathbb F_2^k\to\mathbb F_2^{\bar n}$ of rate
\[\frac k{\bar n}=2^{-c_r\log\log(1/\varepsilon)\sqrt{\log(1/\varepsilon)}}\cdot\varepsilon^2\]
for some universal constant$c_r$, that is encodable in deterministic time$\exp(\exp(\sqrt{\log(1/\varepsilon)}))\cdot k$.\footnote{Following the footnote on the base code, one can get a randomized encoding in time$\operatorname{poly}(1/\varepsilon)\cdot k$ by using a randomized encoding of the based code.} More precisely, there exists a randomized preprocessing step that runs in time$\exp(\sqrt{\log(1/\varepsilon)})\cdot k$ and succeeds with probability$1-2^{-\Omega(k)}$. Once it succeeds, it fixes a deterministic mapping$C$ such that for every$x\in\mathbb F_2^k$, $C(x)$ can be computed in deterministic in the above$O_{\varepsilon}(\bar n)$ time.
If, moreover, for a large enough constant$d$ we are given an explicit family of$\lambda$-spectral$d$-regular3-uniform hypergraphs satisfying$\lambda=\operatorname{poly}(\log d)/\sqrt d$, there is no need for a preprocessing step, and$C$ is explicit.
\paragraph{Theorem36.} Given$k\in\mathbb N$ and$\varepsilon>0$, let$C:\mathbb F_2^k\to\mathbb F_2^{\bar n}$ be the$\varepsilon$-balanced, linear-time encodable code guaranteed to us by Theorem31 with high probability. Then,$C$ also admits the following decoding capabilities.
\begin{enumerate}
\item $C$ is list decodable up to radius$1/2-\beta$ for$\beta=2^{-\frac12\sqrt{\log(1/\varepsilon)}}$ by a randomized algorithm that runs in time$c_1(\varepsilon)\cdot\widetilde O(k)$ and succeeds with probability$1-2^{-\Omega(k)}$.
\item $C$ is uniquely decodable to within distance$(1-\varepsilon)/4$ by a randomized algorithm that runs in time$c_2(\varepsilon)\cdot\widetilde O(k)$ and succeeds with probability$1-2^{-\Omega(k)}$.
\end{enumerate}
Above,$c_1(\varepsilon)=\exp(\exp(\exp(\sqrt{\log(1/\varepsilon)})))$ and$c_2(\varepsilon)=\exp(\exp(\sqrt{\log(1/\varepsilon)}))$.
